% -------------------------------
% Plantilla del TFM en la ESIIAB/ESICR de la UCLM. Versión 0.1 (Preliminar)
%
% Luis de la Ossa
% 
% Compilar con XeLaTeX 
% -------------------------------


% -------------------------------
% Idioma
% -------------------------------
\usepackage[spanish]{babel}


% -------------------------------
% Fuente
% -------------------------------
% Cualquier tamaño de texto: {\fontsize{100pt}{120pt}\selectfont tutexto}
\usepackage{anyfontsize}
% Selección de fuentes.
\usepackage{fontspec}
% Espacio entre líneas
\usepackage{setspace}
% Fuentes especiales
\usepackage{textcomp,marvosym,pifont}


% -------------------------------
% Otros paquetes de uso común
% -------------------------------
% Símbolos matemáticos
\usepackage{amsmath,amsfonts,amssymb}
% Gráficos
\usepackage{graphicx}
% Posición arbitraria de los gráficos
\usepackage[absolute,overlay]{textpos}
% Apéndices
\usepackage{appendix}
% Acrónimos / glosario de siglas
% nohyperlinks evita warnings de hyperref para anchors acro:* que se crean
% en cuerpo después de imprimir el índice (frontmatter).
\usepackage[printonlyused,nohyperlinks]{acronym}
% Árboles de directorios en LaTeX
\usepackage{dirtree}


% -------------------------------
% Esquema de colores
% -------------------------------
% Paquete para la definición de colores
\usepackage[table]{xcolor}
% Esquema de colores
\input{include/colores}


% -------------------------------
% Márgenes
% -------------------------------
% Márgenes de las páginas
\usepackage[
  inner	 =	3cm,   % Margen interior
  outer	 =	2.5cm, % Margen exterior
  top	 =	2.5cm, % Margen superior
  bottom =	2.5cm, % Margen inferior
  includeheadfoot, % Incluye cabecera y pie de página en los márgenes
]{geometry}

% Muestra una regla para comprobar el formato de las páginas (descomentar)
%\usepackage[type=upperleft,showframe,marklength=8mm]{fgruler} 


% -------------------------------
% Interlineado
% -------------------------------
\usepackage{setspace}
\onehalfspacing    % interlineado 1.5
%\doublespacing     % interlineado doble
%\singlespacing     % normal


% -------------------------------
% Renombra tablas y bibliografía
% -------------------------------
\addto\captionsspanish{
	\renewcommand{\listtablename}{Índice de tablas} 
	\renewcommand{\tablename}{Tabla}
	\renewcommand{\bibname}{Referencia bibliográfica}
}


% -------------------------------
% Itemize y enumerate. Control de los espacios.
% -------------------------------
\usepackage{enumitem}
\setitemize{itemsep=0pt}
\setenumerate{itemsep=0pt}
\renewcommand{\labelitemi}{\tiny\bft{$\blacksquare$}}
\renewcommand{\labelitemii}{\bft{$\bullet$}}
\renewcommand{\labelitemiii}{\bft{-}}


% -------------------------------
% Enlaces y colores
% -------------------------------
% hyphens en `url` permite romper URLs largas por guiones y barras, lo que
% evita overfull en la bibliografía (las entradas \url{...} de ref.bib
% suelen exceder el ancho de línea). El entorno sloppypar en TFM.tex
% alrededor de \bibliography{} cierra cualquier resto de overflow residual.
\PassOptionsToPackage{hyphens}{url}
\usepackage{hyperref}
\setlength{\headheight}{15pt}
\emergencystretch=3em
\Urlmuskip=0mu plus 2mu\relax
\emergencystretch=3.5em

% -------------------------------
% Ajuste tipográfico fino y ruptura de identificadores largos
% -------------------------------
% microtype con XeLaTeX solo soporta protrusion; el resto se desactiva
% explícitamente para evitar warnings y errores en cascada.
\usepackage[protrusion=true,expansion=false,
            kerning=false,spacing=false,tracking=false,final]{microtype}

% Romper texto monoespaciado largo (FQN, snake_case) en puntos y guiones bajos
% para evitar overfull boxes con identificadores como
% postgres_demo_service.opendata_demo.gold.movilidad_resumen_municipio.
% Se usa \nolinkurl (de hyperref) porque maneja los caracteres especiales
% _, ., /, - como puntos de ruptura naturales y no crea hipervínculo.
\urlstyle{tt}
\newcommand{\fqn}[1]{{\small\nolinkurl{#1}}}
% \ident: identificador monoespaciado a tamaño normal que SÍ puede cortarse en
% puntos naturales (_ . : - /), para que clases e identificadores largos no
% invadan el margen. Alternativa partible a \texttt para código en línea.
% (\code ya está tomado por el entorno flotante de listados.)
\newcommand{\ident}[1]{{\nolinkurl{#1}}}
%\hypersetup{
%    colorlinks,
%    citecolor=black,
%    filecolor=black,
%    linkcolor=black,
%    urlcolor=black
%}

% -------------------------------
% Gestión de elementos flotantes
% -------------------------------
% Para forzar a que las figuras aparezcan antes del fin de una sección.
% Con la opción [section] se inserta un \FloatBarrier al inicio de cada sección,
% de modo que las figuras flotan pero no cruzan el límite de su sección: tienden
% a colocarse al final de la sección que las referencia, mientras que las tablas
% (con [h]) permanecen junto al texto.
\usepackage[section]{placeins}
% Permite el especificador [H] (colocación exacta, sin flotar) para figuras que
% deben quedar dentro de su subsección y no saltar a la siguiente.
\usepackage{float}

% Gestión de títulos de los elementos flotantes
\usepackage{caption}
% Subfiguras y títulos en subfiguras
\usepackage{subcaption}

% Títulos
\DeclareCaptionFont{tema}{\color{tema}} % Color
\captionsetup{labelfont={tema, bf}}     % Estilo
\captionsetup{font=normalsize}          % Tamaño
\captionsetup{width=.9\linewidth}       % Anchura del título
 % Distancia de la figura al texto (ajustada para control de extensión).
\setlength{\intextsep}{0.45cm}
\setlength{\textfloatsep}{0.45cm}
       

% -------------------------------
% Utilidades para tablas
% -------------------------------
% Para fundir filas en tablas
\usepackage{multirow}
% Para definir columnas con ancho fijo y alineación
\usepackage{array,ragged2e}
\newcolumntype{L}[1]{>{\raggedright\let\newline\\\arraybackslash\hspace{0pt}}m{#1}}
\newcolumntype{C}[1]{>{\centering\let\newline\\\arraybackslash\hspace{0pt}}m{#1}}
\newcolumntype{R}[1]{>{\raggedleft\let\newline\\\arraybackslash\hspace{0pt}}m{#1}}


% -------------------------------
% Para tabular
% -------------------------------
\usepackage{tabto}


% -------------------------------
% Formato de títulos, capítulos, etc.
% -------------------------------
\usepackage{titlesec, titletoc}
% Parte
\titleformat{\part}[display]												     	   % Estilo
	{\Huge \centering \color{tema}}                                                % Formato
	{Parte \thepart }                                                              % Etiqueta
	{0.75ex}                                                                       % Separación
	{\bfseries\fontsize{25pt}{40pt}\selectfont}                                    % Entre etiqueta y título            
	[\thispagestyle{empty}]                                                        % Después

% Capítulo		
\titleformat{\chapter}[block]            			 			                    % Estilo  	
	{\vspace{0cm} \flushright \bfseries \LARGE \color{tema}}  	        		% Formato 
	{\thechapter.}                  			  									% Etiqueta
	{0.75ex}                                     									% Separación
	{}                                           									% Entre etiqueta y título
	[\vspace{-0.25cm} \rule{0.5\textwidth}{1pt}\vspace*{1cm}\thispagestyle{plain}]    % Después      
  
% Sección
\titleformat{\section}
	{\vspace{5pt}  \Large \color{tema}}
	{\thesection .}
	{0.75ex}
	{}

% Subsección
\titleformat{\subsection}
	{\large \color{tema}}
	{\thesubsection .}
	{0.75ex}
	{}
 
% Subsubsección
% Numerada (secnumdepth=3) para que las subsecciones degradadas al fusionar
% capítulos en "Resultados" conserven número y sean referenciables con \ref.
\setcounter{secnumdepth}{3}
\titleformat{\subsubsection}
	{\bf\color{tema}}
	{\thesubsubsection.}
	{0.75ex}
	{}
	 
% Parte (TOC)
\titlecontents{part}[2pc]
  	{\vspace{20pt}\bfseries\filright\Large\color{tema}}
  	{\contentslabel{1.5pc}}{\hspace*{-1.5pc}}
  	{\vspace{5pt}}
  	{}
  
% Capítulo (TOC)
\titlecontents{chapter}[2pc]
  	{\vspace{10pt}\bfseries\large\filright}
  	{\contentslabel{1.5pc}}{\hspace*{-1.5pc}}
  	{\mdseries\titlerule*[0.5pc]{.}\bfseries\contentspage}

% Sección (TOC)
\titlecontents{section}[4pc]
  	{\vspace{5pt}\filright}
  	{\contentslabel{2pc}}{\hspace*{2pc}}
  	{\titlerule*[0.5pc]{.}\contentspage}
  	{\vspace{5pt}}

% Subsección (TOC)
\titlecontents{subsection}[6pc]
  	{\vspace{2.5pt}\filright\small\em}
  	{\contentslabel{3pc}}{\hspace*{-3pc}}
  	{\titlerule*[0.5pc]{.}\contentspage} 
  	
  	
% -------------------------------
% Cabeceras y pie de página
% -------------------------------
\usepackage{fancyhdr}
\pagestyle{fancy}
\fancyhf{}

% Permite incluir la parte. Para ello crea \parttitle
\let\Oldpart\part
\newcommand{\parttitle}{}
\renewcommand{\part}[1]{\Oldpart{#1}\renewcommand{\parttitle}{#1}}

% Cabeceras y pies de página
\fancyhead[LE]{}
\fancyhead[RO]{\slshape \leftmark}
\fancyfoot[LE,RO]{\thepage}
\renewcommand{\footrulewidth}{1pt}
\renewcommand{\headrulewidth}{1pt}
\renewcommand{\chaptermark}[1]{\markboth{#1}{}} % Capítulo (Solo texto)
\renewcommand{\chaptermark}[1]{\markboth{\thechapter. #1}{}} % Capítulo (Número y texto)

% Formato plano para las páginas con títulos (solo incluye el pie de página)
\fancypagestyle{plain}{
\fancyhf{}
\fancyhead{}
\fancyfoot[LE,RO]{\thepage}
\renewcommand{\footrulewidth}{1pt}
\renewcommand{\headrulewidth}{0pt}
}

 
% -------------------------------
% Utilidades
% -------------------------------
% Marcas
\newcommand{\pcite}{\bfr{\large[?]}\,} % Cita pendiente: [?] en rojo y negrita.
\newcommand{\ps}{\bfr{\huge[--}\,} % Pricipio de un bloque en sucio
\newcommand{\fs}{\bfr{\huge--]}\,} % Fin de un bloque en sucio.

% Para generar textos "random"


% -------------------------------
% Algoritmos y código
% -------------------------------

% Se definen entornos específicos para algoritmos y código que permiten
% gestionar los títulos manera similar en ambos casos, y similar al resto
% de elementos flotantes.
\usepackage{newfloat}

\DeclareFloatingEnvironment[ % Algoritmos
  listname = {Índice de algoritmos},
  name = Algoritmo
]{algorithm}

\DeclareFloatingEnvironment[ % Fragmentos de código
  listname = {Índice de listados de código},
  name = Código
]{code}

% Define un estilo para el encabezado de algoritmos y código.
\DeclareCaptionFormat{algcode}{
% OPCIÓN 1 -> Línea debajo del título
%{\color{gris50}\rule{\dimexpr\textwidth\relax}{0.4pt}\vspace{0cm}}
%#1#2#3
%\vspace{-0.2cm}{\color{gris50}\rule{\dimexpr\textwidth\relax}{0.4pt}} % 
% OPCIÓN 2 -> Sin línea debajo del título
{\color{gris50}\rule{\textwidth}{0.4pt}\vspace{0cm}}
#1#2#3
}
% Opciones específicas para algoritmos.
\usepackage[noend]{algpseudocode}
% Tamaño de letra y separadores
\makeatletter
\renewcommand{\ALG@beginalgorithmic}{{\color{gris50}\small\hrule\vskip10pt}}
\makeatother
% Variables y funciones
\newcommand{\va}[1]{{\it {#1}}}				
\newcommand{\fu}[1]{{\textsc{#1}}}          
% Palabras clave
\renewcommand{\algorithmicrepeat}{\bft{repetir}}
\renewcommand{\algorithmicuntil}{\bft{hasta}}
\renewcommand{\algorithmicend}{\bft{fin}}
\renewcommand{\algorithmicif}{\bft{si}}
\renewcommand{\algorithmicthen}{\bft{entonces}}
\renewcommand{\algorithmicelse}{\bft{si no}}
\newcommand{\algorithmicelsif}{\algorithmicelse\ \algorithmicif}
\newcommand{\algorithmicendif}{\algorithmicend\ \algorithmicif}
\renewcommand{\algorithmicfor}{\bft{para}}
\renewcommand{\algorithmicforall}{\bft{para cada}}
\renewcommand{\algorithmicwhile}{\bft{mientras}}
\renewcommand{\algorithmicdo}{\bft{hacer}}
\renewcommand{\algorithmicprocedure}{\bft{procedimiento}}
\renewcommand{\algorithmicreturn}{\bft{devolver}}
\renewcommand{\algorithmicfunction}{\bft{función}}
\renewcommand{\algorithmicloop}{\bft{iterar}}

% Opciones específicas para código.
\usepackage{listings}

% --- Lenguajes personalizados para los listados del TFM ---
% PowerShell (scripts de orquestación e infraestructura)
\lstdefinelanguage{PowerShell}{
    morekeywords={param, Param, function, Function, if, else, elseif, foreach, ForEach-Object, for, while, do, switch, return, break, continue, try, catch, finally, throw, in, begin, process, end, Write-Host, Write-Output, Write-Error, Write-Warning, Get-ChildItem, Set-Location, New-Item, Remove-Item, Test-Path, Import-Module, Export-ModuleMember, Invoke-Expression, Invoke-RestMethod, Invoke-WebRequest, Start-Process, Stop-Process, kubectl, helm, kind, python, pnpm},
    sensitive=true,
    morecomment=[l]{\#},
    morecomment=[s]{<\#}{\#>},
    morestring=[b]",
    morestring=[b]'
}

% TypeScript (consola web Next.js)
\lstdefinelanguage{TypeScript}{
    morekeywords={abstract, any, as, async, await, boolean, break, case, catch, class, const, continue, declare, default, delete, do, else, enum, export, extends, false, finally, for, from, function, get, if, implements, import, in, instanceof, interface, let, module, new, null, number, of, package, private, protected, public, readonly, return, set, static, string, super, switch, symbol, this, throw, true, try, type, typeof, undefined, var, void, while, with, yield},
    sensitive=true,
    morecomment=[l]{//},
    morecomment=[s]{/*}{*/},
    morestring=[b]",
    morestring=[b]',
    morestring=[b]`
}

% YAML (configuración operativa)
\lstdefinelanguage{yaml}{
    keywords={true, false, null, yes, no, on, off},
    sensitive=false,
    morecomment=[l]{\#},
    morestring=[b]",
    morestring=[b]',
    otherkeywords={:, -, >, |, !}
}

% Turtle / SHACL (RDF en formato Terse Triple Language)
\lstdefinelanguage{Turtle}{
    morekeywords={@prefix, @base, a, true, false, sh:NodeShape, sh:PropertyShape, sh:targetClass, sh:property, sh:path, sh:minCount, sh:maxCount, sh:datatype, sh:class, sh:nodeKind, sh:severity, sh:IRI, sh:Literal, sh:Violation, sh:Warning, sh:Info, sh:in, dcat:Catalog, dcat:Dataset, dcat:Distribution, dcat:DataService, dcterms:title, dcterms:description},
    sensitive=true,
    morecomment=[l]{\#},
    morestring=[b]",
    morestring=[b]'
}

% JSON-LD (serialización del catálogo)
\lstdefinelanguage{jsonld}{
    morekeywords={true, false, null},
    sensitive=true,
    morecomment=[l]{//},
    morestring=[b]"
}

% Bash auxiliar (cuando aplique)
\lstdefinelanguage{bash}{
    morekeywords={if, then, else, elif, fi, for, do, done, while, case, esac, function, return, in, export, source, echo, cd, ls, mkdir, rm, cp, mv, kubectl, helm, kind, python, pnpm, curl},
    sensitive=true,
    morecomment=[l]{\#},
    morestring=[b]",
    morestring=[b]'
}

% --- Aspecto común de todos los listados ---
% Estilo visual: caja con fondo gris claro, banda superior coloreada, numeración
% de líneas tenue. Pensado para alta legibilidad en impresión y pantalla.
\definecolor{listingbg}{RGB}{246,246,248}
\definecolor{listingrule}{RGB}{210,210,215}
\lstset{
    language=Python,
    keywordstyle=\bfseries\color{tema},
    commentstyle=\itshape\color{gris50},
    stringstyle=\color{tema!70!black},
    identifierstyle=\color{black},
    breaklines=true,
    breakatwhitespace=false,
    postbreak=\mbox{\textcolor{gris50}{$\hookrightarrow$}\space},
    breakindent=10pt,
    showstringspaces=false,
    basicstyle=\footnotesize\ttfamily,
    tabsize=2,
    captionpos=b,
    numbers=left,
    numberstyle=\tiny\color{gris50},
    numbersep=8pt,
    backgroundcolor=\color{listingbg},
    rulecolor=\color{listingrule},
    framerule=0.5pt,
    framesep=6pt,
    xleftmargin=12pt,
    xrightmargin=4pt,
    aboveskip=0.5em,
    belowskip=0.5em,
    frame=single,
    columns=fullflexible,
    keepspaces=true,
    literate={á}{{\'a}}1 {é}{{\'e}}1 {í}{{\'i}}1 {ó}{{\'o}}1 {ú}{{\'u}}1
             {Á}{{\'A}}1 {É}{{\'E}}1 {Í}{{\'I}}1 {Ó}{{\'O}}1 {Ú}{{\'U}}1
             {ñ}{{\~n}}1 {Ñ}{{\~N}}1 {¿}{{?`}}1 {¡}{{!`}}1
}

% -------------------------------
% Cacacteres pifont de uso común (con color)
% -------------------------------
\newcommand{\vmarkt}{\textcolor{tema}{\ding{52}} }  % Checked (V) Tema
\newcommand{\vmarkr}{\textcolor{rojo}{\ding{52}} }  % Checked (V) Rojo
\newcommand{\vmarka}{\textcolor{azul}{\ding{52}} }  % Checked (V) Azul

\newcommand{\xmarkt}{\textcolor{tema}{\ding{55}} }  % Checked (X) Tema
\newcommand{\xmarkr}{\textcolor{rojo}{\ding{55}} }  % Checked (X) Rojo
\newcommand{\xmarka}{\textcolor{azul}{\ding{55}} }  % Checked (X) Azul

\newcommand{\lmarkt}{\textcolor{tema}{\ding{46}} }  % Lápiz Tema
\newcommand{\lmarkr}{\textcolor{rojo}{\ding{46}} }  % Lápiz Rojo
\newcommand{\lmarka}{\textcolor{azul}{\ding{46}} }  % Lápiz Azul

\newcommand{\hmarkt}{\textcolor{tema}{\ding{43}} }  % Mano Tema
\newcommand{\hmarkr}{\textcolor{rojo}{\ding{43}} }  % Mano Roja 
\newcommand{\hmarka}{\textcolor{azul}{\ding{43}} }  % Mano Azula 

\newcommand{\fmarkt}{\textcolor{tema}{\ding{220}} }  % Flechza Tema
\newcommand{\fmarkr}{\textcolor{rojo}{\ding{220}} }  % Flecha Roja 
\newcommand{\fmarka}{\textcolor{azul}{\ding{220}} }  % Flecha Azul

% -------------------------------
% Listados y figuras con captura opcional
% -------------------------------
\renewcommand{\lstlistingname}{Código}
\renewcommand{\lstlistlistingname}{Índice de Códigos}

% Figuras de la memoria: ancho por defecto 0.8\linewidth y altura limitada a
% 0.45\textheight para que las figuras (sobre todo los diagramas) no ocupen
% demasiado y ayudar a ajustar el cuerpo al límite de 50 páginas; conserva
% relación de aspecto. Si la figura no existe, muestra "Captura pendiente".
\newcommand{\figuraMemoria}[4][0.8\linewidth]{%
\begin{figure}[htbp]
\centering
\IfFileExists{figs/#2}{%
  \includegraphics[width=#1,height=0.45\textheight,keepaspectratio]{figs/#2}%
}{%
  \fbox{%
    \begin{minipage}[c][5cm][c]{0.85\linewidth}
    \centering
    \textbf{Captura pendiente}\\[0.4cm]
    {\ttfamily\detokenize{#2}}
    \end{minipage}%
  }%
}
\caption{#3}
\label{#4}
\end{figure}%
}
